% GENERATED FILE. Edit canonical lessons, then run npm run generate:books.
% canonical-source-hash: fnv1a64:868957643f51d0d3
% canonical-lessons: ES-C402-costar, ES-C402-camisa, ES-C402-camiseta, ES-C402-talla, ES-C402-gorro, ES-C402-shopping-recall-1, ES-C402-shopping-recall-2

\chapter{Five Words in a Clothes Shop}
\label{ch:camisa-camiseta-talla-gorro-costar}
\hlchaptermodality{402}

\begin{chapteropening}
\textbf{By the end of this chapter:} \emph{I can name a shirt, T-shirt, size and cap, and ask or say what something costs.}

Complete ES-C402-gorro: identify four clothes-shop words and ask what an item costs.
\end{chapteropening}

\section[costar]{costar — when the price stands settled}

\label{lesson:ES-C402-costar}



\begin{quote}
Say \emph{the euro} and \emph{the price}. (\emph{El euro, el precio}.)
\end{quote}

\subsection*{What to know first}
\begin{itemize}
\raggedright
  \item el precio — the number this verb asks about.
  \item el euro — one unit the answer can use.
\end{itemize}

\begin{sounds}
\begin{itemize}
\raggedright
  \item \emph{cos-TAR}: keep the first \emph{s} clear. \textbf{Costar} differs from \textbf{contar} only in that middle consonant: \emph{cos-tar}, not \emph{con-tar}.
\end{itemize}
\end{sounds}

\begin{cousinweb}
\textbf{costar} — \textquotedblleft{}to cost.\textquotedblright{}

It descends from Latin \textbf{constare}, \textquotedblleft{}to stand together\textquotedblright{} and then \textquotedblleft{}to be settled or established.\textquotedblright{} A settled price is what an item \textbf{costs}. English \textbf{cost} reached the same Latin family through French.

Use the third-person question for one item: \textbf{¿Cuánto cuesta?} — \textquotedblleft{}How much does it cost?\textquotedblright{} The answer can be \textbf{Cuesta diez euros}.
\end{cousinweb}

\subsection*{Your turn}
Say these aloud:

\begin{itemize}
\raggedright
  \item \emph{costar}
  \item \emph{¿Cuánto cuesta?}
  \item \emph{Cuesta diez euros.}
\end{itemize}

\begin{tcolorbox}[breakable,colback=teal!4,colframe=teal!35!black,title={Before you move on}]
Ask how much it costs. (\emph{¿Cuánto cuesta?})

Source: \href{https://dle.rae.es/costar}{Diccionario de la lengua española: costar}.
\end{tcolorbox}
\section[la camisa]{la camisa — a garment with a long family}

\label{lesson:ES-C402-camisa}



\begin{quote}
Ask how much it costs. (\emph{¿Cuánto cuesta?})
\end{quote}

\begin{sounds}
\begin{itemize}
\raggedright
  \item \emph{ca-MI-sa}: three open syllables, with the weight on \emph{mi}.
\end{itemize}
\end{sounds}

\begin{cousinweb}
\textbf{la camisa} — \textquotedblleft{}the shirt.\textquotedblright{}

Spanish continues Late Latin \textbf{camisia}, a garment worn close to the body. English \textbf{chemise} and \textbf{camisole} travelled by another route from the same old garment word, so their opening sound changed while Spanish kept \emph{camisa}.

In the shop: \textbf{¿Cuánto cuesta la camisa?}
\end{cousinweb}

\subsection*{Your turn}
Say these aloud:

\begin{itemize}
\raggedright
  \item \emph{la camisa}
  \item \emph{¿Cuánto cuesta la camisa?}
  \item \emph{La camisa cuesta diez euros.}
\end{itemize}

\begin{tcolorbox}[breakable,colback=teal!4,colframe=teal!35!black,title={Before you move on}]
Name the shirt, then ask its cost. (\emph{La camisa. ¿Cuánto cuesta?})

Source: \href{https://dle.rae.es/camisa}{Diccionario de la lengua española: camisa}.
\end{tcolorbox}
\section[la camiseta]{la camiseta — the smaller shirt}

\label{lesson:ES-C402-camiseta}



\begin{quote}
Say \emph{the shirt} and ask what it costs. (\emph{La camisa. ¿Cuánto cuesta?})
\end{quote}

\begin{sounds}
\begin{itemize}
\raggedright
  \item \emph{ca-mi-SE-ta}: keep all four vowels; the stress moves to \emph{se}.
\end{itemize}
\end{sounds}

\begin{cousinweb}
\textbf{la camiseta} — \textquotedblleft{}the T-shirt.\textquotedblright{}

The build is visible: \textbf{camisa} plus \textbf{-eta}, a suffix that can make something smaller. The lighter, smaller shirt became the ordinary T-shirt. Keep the two shop labels distinct: \textbf{camisa} for a shirt, \textbf{camiseta} for a T-shirt.
\end{cousinweb}

\subsection*{Your turn}
Say these aloud:

\begin{itemize}
\raggedright
  \item \emph{la camisa}
  \item \emph{la camiseta}
  \item \emph{¿Cuánto cuesta la camiseta?}
\end{itemize}

\begin{tcolorbox}[breakable,colback=teal!4,colframe=teal!35!black,title={Before you move on}]
Give the two garment words, larger category first. (\emph{La camisa, la camiseta}.)

Source: \href{https://dle.rae.es/camiseta}{Diccionario de la lengua española: camiseta}.
\end{tcolorbox}
\section[la talla]{la talla — the cut that became a size}

\label{lesson:ES-C402-talla}



\begin{quote}
Name a shirt and a T-shirt. (\emph{Una camisa, una camiseta}.)
\end{quote}

\begin{sounds}
\begin{itemize}
\raggedright
  \item \emph{TA-ya} for most speakers: the written \textbf{ll} commonly has the same sound as \textbf{y}. $\to$ pronunciation reference
\end{itemize}
\end{sounds}

\begin{cousinweb}
\textbf{la talla} — \textquotedblleft{}the size,\textquotedblright{} especially a clothing size.

The noun belongs to \textbf{tallar}, \textquotedblleft{}to cut or shape.\textquotedblright{} A garment is cut to a measure, and that cut becomes its \textbf{talla}. In a shop, \textbf{¿Qué talla?} asks \textquotedblleft{}What size?\textquotedblright{}
\end{cousinweb}

\subsection*{Your turn}
Say these aloud:

\begin{itemize}
\raggedright
  \item \emph{la talla}
  \item \emph{¿Qué talla?}
  \item \emph{¿Qué talla tiene la camiseta?}
  \item \emph{¿Cuánto cuesta la camisa?}
\end{itemize}

\begin{tcolorbox}[breakable,colback=teal!4,colframe=teal!35!black,title={Before you move on}]
Ask \textquotedblleft{}What size?\textquotedblright{} (\emph{¿Qué talla?})

Source: \href{https://dle.rae.es/talla}{Diccionario de la lengua española: talla}.
\end{tcolorbox}
\section[el gorro]{el gorro — the cap with no invented origin}

\label{lesson:ES-C402-gorro}



\begin{quote}
Ask the size of the T-shirt, then ask its cost. (\emph{¿Qué talla tiene la camiseta? ¿Cuánto cuesta?})
\end{quote}

\begin{sounds}
\begin{itemize}
\raggedright
  \item \emph{GO-rro}: hold the rolled \textbf{rr} clear. Contrast \textbf{gorro}, with \emph{rr}, and \textbf{gordo}, with \emph{rd}. The new word does not contain a \emph{d}.
\end{itemize}
\end{sounds}

\begin{cousinweb}
\textbf{el gorro} — \textquotedblleft{}the brimless cap,\textquotedblright{} often round or close-fitting.

There is no secure etymology to turn into a story here. The Academy marks the origin of \textbf{gorra} as uncertain and derives \textbf{gorro} from it. Learn the shape instead: a \textbf{gorra} normally has a peak; a \textbf{gorro} is round or brimless. They are not masculine and feminine forms of one object.
\end{cousinweb}

\subsection*{Your turn}
Say these aloud:

\begin{itemize}
\raggedright
  \item \emph{el gorro}
  \item \emph{gorro, no gordo}
  \item \emph{¿Qué talla tiene el gorro?}
  \item \emph{¿Cuánto cuesta el gorro?}
\end{itemize}

\begin{tcolorbox}[breakable,colback=teal!4,colframe=teal!35!black,title={Before you move on}]
Name the brimless cap and ask its cost. (\emph{El gorro. ¿Cuánto cuesta?})

Sources: \href{https://dle.rae.es/gorro}{DLE: gorro} and \href{https://dle.rae.es/gorra}{DLE: gorra}.
\end{tcolorbox}
\section[(review)]{shopping recall one — rebuild the counter exchange}

\label{lesson:ES-C402-shopping-recall-1}



\begin{quote}
Give the five-word chain. (\emph{Costar, la camisa, la camiseta, la talla, el gorro}.)
\end{quote}

\begin{grammarlens}[title={What you've built}]
You can distinguish \textbf{camisa} from \textbf{camiseta}, ask \textbf{¿Qué talla?}, name \textbf{el gorro}, and ask \textbf{¿Cuánto cuesta?}
\end{grammarlens}

\subsection*{Your turn}
Say these aloud:

\begin{itemize}
\raggedright
  \item \emph{la camisa y la camiseta}
  \item \emph{¿Qué talla?}
  \item \emph{el gorro}
  \item \emph{¿Cuánto cuesta?}
\end{itemize}

\begin{tcolorbox}[breakable,colback=teal!4,colframe=teal!35!black,title={Before you move on}]
Ask the size and cost of the T-shirt. (\emph{¿Qué talla tiene la camiseta? ¿Cuánto cuesta?})
\end{tcolorbox}
\section[(review)]{shopping recall two — choose, size, and price}

\label{lesson:ES-C402-shopping-recall-2}



\begin{quote}
Without looking back, name the four shop nouns. (\emph{La camisa, la camiseta, la talla, el gorro}.)
\end{quote}

\subsection*{Your turn}
Say these aloud:

\begin{itemize}
\raggedright
  \item \emph{Busco una camiseta.}
  \item \emph{¿Qué talla?}
  \item \emph{También busco un gorro.}
  \item \emph{¿Cuánto cuesta?}
\end{itemize}

\begin{tcolorbox}[breakable,colback=teal!4,colframe=teal!35!black,title={Before you move on}]
Build a two-question shop exchange using a garment, its size, and its price. (\emph{Busco una camiseta. ¿Qué talla? ¿Cuánto cuesta?})
\end{tcolorbox}
